% Môn: Toán
% Lớp: 10
% Chương: Chương I. Mệnh đề và tập hợp
% Bài: Bài 2. Tập hợp và các phép toán trên tập hợp · Phép giao, hợp, hiệu và phần bù của hai tập hợp
% Số câu: 56
% App đọc file này trực tiếp — sửa rồi Commit trên GitHub, bấm Đồng bộ GitHub .tex

% Bài 2. Tập hợp và các phép toán trên tập hợp



% ===== Câu 2 =====
% ID: T10C1B2-01-TL
% Mức: NB
\dangbt{Phép giao, hợp, hiệu và phần bù của hai tập hợp}
\begin{bt}
% ID: T10C1B2-01-TL
% Mức: NB
Cho hai tập hợp $A=\{n \in \mathbb{Z} \mid -2 \lt  n \leq 4\}$ và $B=\{-1 ; 0;1;5\}$. Hãy xác định tập hợp $A \cup B$, $A \setminus B$.
\loigiai{
Ta có $A=\{-1; 0 ; 1 ; 2 ; 3 ; 4\}$, $B=\{-1 ; 0;1;5\}$.
Suy ra $A \cup B=\{-1; 0 ; 1 ; 2 ; 3 ; 4; 5\}$ và $A \setminus B=\{2 ; 3 ; 4\}$.
}
\end{bt}

% ===== Câu 5 =====
% ID: T10C1B2-02-DS
% Mức: VD
\begin{ex}
% ID: T10C1B2-02-DS
% Mức: VD
Cho các tập hợp $A=\{-3 ;-2 ;-1 ; 0 ; 1 ; 2 ; 3\}$, $B=\{0 ; 1 ; 4 ; 5\}$ và $C=\{-4 ;-3 ; 1 ; 2 ; 5 ; 6\}$.
\choiceTF
{$A\cup B=\{-3 ;-2 ;-1 ; 0 ; 1 ; 2 ; 3 ; 4 ; 5\}$}
{$A\cap B=\{0\}$}
{\True $(A\cup B)\cap C=\{-3 ; 1 ; 2 ; 5\}$}
{$A\cap B\cap C=\{1\}$}
\loigiai{
\begin{itemchoice}
\item (Sai) $A\cup B=\{-3 ;-2 ;-1 ; 0 ; 1 ; 2 ; 3 ; 4 ; 5\}$. (Vì): \cup B)$ và $C$ là tập hợp chứa các phần tử chung của hai tập hợp này. Ta có $
\item (Sai) $A\cap B=\{0\}$. (Vì): Giao của hai tập hợp $A$ và $B$ là tập hợp chứa tất cả các phần tử vừa thuộc $A$ vừa thuộc $B$. Ta có $A\cap B=\{-3 ;-2 ;-1 ; 0 ; 1 ; 2 ; 3\} \cap \{0 ; 1 ; 4 ; 5\} = \{0 ; 1\}$. Mệnh đề cho là $\{0\}$, nên sai
\item (Đúng) $(A\cup B)\cap C=\{-3 ; 1 ; 2 ; 5\}$. (Vì): Hợp của $A$ và $B$ là $A\cup B=\{-3 ;-2 ;-1 ; 0 ; 1 ; 2 ; 3 ; 4 ; 5\}$. Giao của $
\item (Sai) $A\cap B\cap C=\{1\}$.
\end{itemchoice}
}
\end{ex}

% ===== Câu 12 =====
% ID: T10C1B2-03-TN
% Mức: NB
\begin{ex}
% ID: T10C1B2-03-TN
% Mức: NB
Cho tập hợp $A=\{x \in \mathbb{R} \mid 2 \leqslant x<5\}$. Phần bù của tập hợp $A$ trong $\mathbb{R}$ là tập nào sau đây?
\choice
{$[5 ;+\infty)$}
{$(-\infty ; 2)$}
{$(-\infty ; 2] \cup(5 ;+\infty)$}
{\True $(-\infty ; 2) \cup[5 ;+\infty)$}
\loigiai{
Ta có $A = \{x \in \mathbb{R} \mid 2 \leqslant x<5\}
 = [2;5)$.
Khi đó phần bù của tập hợp $A$ trong $\mathbb{R}$ là
 $\mathrm{C_{\mathbb{R}}A}
 = \mathbb{R} \setminus A = (-\infty ; 2) \cup[5 ;+\infty)$.
}
\end{ex}

% ===== Câu 22 =====
% ID: T10C1B2-06-TN
% Mức: NB
\begin{ex}
% ID: T10C1B2-06-TN
% Mức: NB
Cho hai tập hợp $A=\{-1;2;3;5;7\}$, $B=\{1;2;3;4;5\}$. Khi đó giao của hai tập hợp là
\choice
{\True $\{2;3;5\}$}
{$\{-1;2;3\}$}
{$\{1;2;3;4;5\}$}
{$\{2;3;4;5\}$}
\loigiai{
Giao của hai tập hợp $A$ và $B$, ký hiệu là $A\cap B$, là tập hợp gồm các phần tử vừa thuộc $A$ vừa thuộc $B$. 
Ta có $A=\{-1;2;3;5;7\}$ và $B=\{1;2;3;4;5\}$. 
Các phần tử chung của $A$ và $B$ là $2, 3, 5$. 
Vậy $A\cap B = \{2;3;5\}$.
}
\end{ex}

% ===== Câu 26 =====
% ID: T10C1B2-08-TN
% Mức: NB
\begin{ex}
% ID: T10C1B2-08-TN
% Mức: NB
Cho tập $A=\{0 ; 2 ; 4 ; 6 ; 8\}$; $B=\{3 ; 4 ; 5 ; 6 ; 7\}$. Tập $A \setminus B$ là
\choice
{$\{0 ; 6 ; 8\}$}
{\True $\{0 ; 2 ; 8\}$}
{$\{3 ; 6 ; 7\}$}
{$\{0 ; 2\}$}
\loigiai{
Vì $A \setminus B$ gồm các phần tử thuộc tập hợp $A$ và không thuộc $B$ nên $A \setminus B=\{0 ; 2 ; 8\}$.
}
\end{ex}

% ===== Câu 27 =====
% ID: T10C1B2-09-TN
% Mức: NB
\begin{ex}
% ID: T10C1B2-09-TN
% Mức: NB
Cho hai tập hợp $A=\{1;2;3;4;5;6;7\}$, $B=\{5;6;7;8\}$. Tìm tập $C=A\cup B$.
\choice
{\True $C=\{1;2;3;4;5;6;7;8\}$}
{$C=\{1;2;3;4;8\}$}
{$C=\{5;6;7\}$}
{$C=\{1;2;3;4\}$}
\loigiai{
Hợp của hai tập hợp $A$ và $B$, ký hiệu là $A \cup B$, là tập hợp gồm tất cả các phần tử thuộc $A$ hoặc thuộc $B$ hoặc thuộc cả hai tập hợp. 
Ta có $A = \{1; 2; 3; 4; 5; 6; 7\}$ và $B = \{5; 6; 7; 8\}$. 
Khi đó, $A \cup B = \{1; 2; 3; 4; 5; 6; 7\} \cup \{5; 6; 7; 8\} = \{1; 2; 3; 4; 5; 6; 7; 8\}$. 
Vậy tập $C = A \cup B = \{1; 2; 3; 4; 5; 6; 7; 8\}$.
}
\end{ex}

% ===== Câu 32 =====
% ID: T10C1B2-14-TN
% Mức: NB
\begin{ex}
% ID: T10C1B2-14-TN
% Mức: NB
Cho $A=\{0;1;2\}$, $B=\{-1;0;1\}$. Khi đó $A\cap B$ là
\choice
{$\{0\}$}
{$\{-1;0;1\}$}
{\True $\{0;1\}$}
{$\{1\}$}
\loigiai{
Phần tử chung của hai tập hợp $A$ và $B$ là 0 và 1. Do đó, $A\cap B = \{0;1\}$.
}
\end{ex}

% ===== Câu 36 =====
% ID: T10C1B2-18-TN
% Mức: NB
\begin{ex}
% ID: T10C1B2-18-TN
% Mức: NB
Cho tập hợp $A=[3;8)$, $B=(-\infty;4)$. Tìm $A\cap B$
\choice
{\True [3;4)}
{[3;4]}
{(3;4)}
{(3;4]}
\loigiai{
Tập hợp $A$ bao gồm các số thực lớn hơn hoặc bằng 3 và nhỏ hơn 8. Tập hợp $B$ bao gồm các số thực nhỏ hơn 4. Giao của hai tập hợp $A$ và $B$ là tập hợp các phần tử thuộc cả $A$ và $B$. Do đó, $A\cap B$ bao gồm các số thực lớn hơn hoặc bằng 3 và nhỏ hơn 4. Vậy $A\cap B=[3;4)$.
}
\end{ex}

% ===== Câu 37 =====
% ID: T10C1B2-19-TN
% Mức: NB
\begin{ex}
% ID: T10C1B2-19-TN
% Mức: NB
Tập hợp $X=(-\infty;2]\cap (-6;+\infty)$ là
\choice
{$(-6;2)$}
{\True $(-6;2]$}
{$[-\infty;2]$}
{$(-6;+\infty)$}
\loigiai{
Ta có $X=(-\infty;2]\cap (-6;+\infty)$. Giao của hai tập hợp này là tập hợp các phần tử thuộc cả hai tập hợp. Do đó, $X=(-6;2]$.
}
\end{ex}

% ===== Câu 40 =====
% ID: T10C1B2-22-TN
% Mức: NB
\begin{ex}
% ID: T10C1B2-22-TN
% Mức: NB
Cho $A=[-1 ; 3]$; $B=(2 ; 5)$. Tìm mệnh đề $\textbf{sai}$.
\choice
{$B \setminus A=(3 ; 5)$}
{$A \cap B=(2 ; 3]$}
{$A \setminus B=[-1 ; 2]$}
{\True $A \cup B=[-1 ; 5]$}
\loigiai{
Ta có $A \cup B=[-1 ; 5)$.
}
\end{ex}

% ===== Câu 41 =====
% ID: T10C1B2-23-TN
% Mức: TH
\begin{ex}
% ID: T10C1B2-23-TN
% Mức: TH
Cho $A\subset B$. Khẳng định nào sau đây là đúng?
\choice
{$A \cup B = A$}
{\True $A \cap B = A$}
{$A \setminus B = \emptyset$}
{$B \setminus A = B$}
\loigiai{
Vì $A\subset B$ nên mọi phần tử của $A$ đều thuộc $B$. Do đó, tập hợp các phần tử thuộc cả $A$ và $B$ chính là tập $A$. Vậy $A \cap B = A$.
}
\end{ex}

% ===== Câu 49 =====
% ID: T10C1B2-31-TN
% Mức: VD
\begin{ex}
% ID: T10C1B2-31-TN
% Mức: VD
Trong các tập hợp sau, tập hợp nào là tập hợp rỗng?
\choice
{$A=\{x \in \mathbb{Z}|~ x<1\}$}
{$B=\left\{x \in \mathbb{Z} \mid 6 x^2-7 x+1=0\right\}$}
{\True $C=\left\{x \in \mathbb{Q} \mid x^2-4 x+2=0\right\}$}
{$D=\left\{x \in \mathbb{R} \mid x^2-4 x+3=0\right\}$}
\loigiai{
Xét tập hợp $C = \{x \in \mathbb{Q} \mid x^2 - 4x + 2 = 0\}$. Phương trình $x^2 - 4x + 2 = 0$ có hai nghiệm là $x = 2 + \sqrt{2}$ và $x = 2 - \sqrt{2}$. Cả hai nghiệm này đều không thuộc tập số hữu tỉ $\mathbb{Q}$. Do đó, tập hợp $C$ không chứa phần tử nào, tức là $C = \varnothing$.
Xét tập hợp $A = \{x \in \mathbb{Z} \mid x \lt  1\}$. Tập hợp này chứa các số nguyên nhỏ hơn 1, ví dụ như $0, -1, -2, \dots$. Do đó, $A \neq \varnothing$.
Xét tập hợp $B = \{x \in \mathbb{Z} \mid 6x^2 - 7x + 1 = 0\}$. Phương trình $6x^2 - 7x + 1 = 0$ có hai nghiệm là $x = 1$ và $x = \frac{1}{6}$. Trong hai nghiệm này, chỉ có $x = 1$ thuộc tập số nguyên $\mathbb{Z}$. Do đó, $B = \{1\}$, suy ra $B \neq \varnothing$.
Xét tập hợp $D = \{x \in \mathbb{R} \mid x^2 - 4x + 3 = 0\}$. Phương trình $x^2 - 4x + 3 = 0$ có hai nghiệm là $x = 1$ và $x = 3$. Cả hai nghiệm này đều thuộc tập số thực $\mathbb{R}$. Do đó, $D = \{1; 3\}$, suy ra $D \neq \varnothing$.
Vậy, tập hợp rỗng là tập hợp $C$.
}
\end{ex}

% ===== Câu 50 =====
% ID: T10C1B2-32-TN
% Mức: VD
\begin{ex}
% ID: T10C1B2-32-TN
% Mức: VD
Trong các tập sau đây, tập hợp nào có đúng hai tập hợp con?
\choice
{$\{x, y\}$}
{\True $\{x\}$}
{$\{0, x\}$}
{$\{0; x, y\}$}
\loigiai{
$\{x, y\}$ có $2^2=4$ tập con.
$\{x\}$ có $2^1=2$ tập con là $\{x\}$ và $\varnothing$.
$\{0, x\}$ có $2^2=4$ tập con.
$\{0; x, y\}$ có $z^3=8$ tập con.
}
\end{ex}
